% Authoritative exercise source. Preserve all derivations and diagrams.

\ExerciseStart{G3.11}{Verificare una porta composta}{Collegamento}

\BookHeading{Problema, dati e obiettivo}

Dimostrare che \((I\otimes H)\,CZ\,(I\otimes H)\) realizza CNOT con primo qubit di controllo. Verificare l’intera base e la preparazione \((\ket{00}+i\ket{10})/\sqrt2\).

\ExTheory{\ExRef{c3-sintesi}{3.10}}{La porta \(CZ=\operatorname{diag}(1,1,1,-1)\) applica Z al secondo qubit soltanto se il primo vale uno. L’identità \(HZH=X\) permette di convertire una fase controllata in un’inversione controllata.}

\ExPrediction{Nel sottospazio con controllo zero le due H dovrebbero cancellarsi. Nel sottospazio con controllo uno la fase cambia la ricombinazione delle componenti del bersaglio.}

\ExFigure{\begin{tikzpicture}[>=Stealth,exlabel]
\node[exbox] (n0) at (0.0,0.0) {$I\otimes H$};\node[exbox] (n1) at (3.6,0.0) {$CZ$};\draw[exother] (n0)--(n1);\node[exbox] (n2) at (7.2,0.0) {$I\otimes H$};\draw[exother] (n1)--(n2);
\end{tikzpicture}}{Le tre operazioni compongono una sola porta logica controllata.}


\begin{ExTasks}

\item Calcolare ZH e poi HZH.

\item Esaminare il sottospazio con controllo zero.

\item Esaminare quello con controllo uno.

\item Scrivere le quattro immagini di base e la matrice complessiva.

\item Applicare la linearità alla preparazione assegnata.

\end{ExTasks}

\ExSolution

\ExStep{Calcolare il prodotto sul bersaglio}

\ExEq{ZH&=\tfrac1{\sqrt2}\begin{pmatrix}1&1\\-1&1\end{pmatrix},\\
HZH&=\tfrac12\begin{pmatrix}1-1&1+1\\1+1&1-1\end{pmatrix}
=\begin{pmatrix}0&1\\1&0\end{pmatrix}=X.}

\ExStep{Controllo zero}

CZ agisce come I sul bersaglio. Il prodotto è \(HIH=H^2=I\), quindi \(00\mapsto00\) e \(01\mapsto01\).

\ExStep{Controllo uno}

Il prodotto sul bersaglio è HZH=X. Quindi \(10\mapsto11\) e \(11\mapsto10\). Per esempio,
\ExEq{\ket{10}&\mapsto(\ket{10}+\ket{11})/\sqrt2\\
&\mapsto(\ket{10}-\ket{11})/\sqrt2\mapsto\ket{11}.}

\ExStep{Ricomporre la matrice}

\ExEq{(I\otimes H)CZ(I\otimes H)
&=\begin{pmatrix}1&0&0&0\\0&1&0&0\\0&0&0&1\\0&0&1&0\end{pmatrix}.}
La verifica sui quattro vettori di base determina l’azione su tutto lo spazio.

\ExStep{Applicare la sovrapposizione}

\ExEq{C\frac{\ket{00}+i\ket{10}}{\sqrt2}
&=\frac{C\ket{00}+iC\ket{10}}{\sqrt2}
=\frac{\ket{00}+i\ket{11}}{\sqrt2}.}
La norma è uno e la fase relativa \(i\) resta nel risultato. Le probabilità sono \(1/2\) sui due rami, senza che sia stato misurato il controllo.

\par\noindent\begin{minipage}{\linewidth}
\ExFigure{\begin{tikzpicture}[>=Stealth,exlabel]
\node[exbox] (n0) at (0.0,0.0) {controllo $0$: $H^2$};\node[exbox] (n1) at (3.6,0.0) {bersaglio invariato};\draw[exother] (n0)--(n1);\node[exbox] (n2) at (7.2,0.0) {controllo $1$: $HZH=X$};\draw[exother] (n1)--(n2);\node[exbox] (n3) at (0.0,-1.6) {bersaglio invertito};\draw[exother] (n2.south)--++(0,-.4)-|(n3.north);
\end{tikzpicture}}{Decomposizione nei due sottospazi ortogonali del controllo.}


\ExInterpret{L’identità di matrici certifica la decomposizione anche sulle sovrapposizioni. Verificare un solo ingresso avrebbe controllato soltanto una colonna della porta.}

\ExLab{Nella sezione 3.10 confrontare «Decomposizione su |10\(\rangle\)» e «Decomposizione su una sovrapposizione». Seguire i tre stati intermedi e verificare che la fase del ramo inizialmente immaginario sia conservata.}


\end{minipage}
